\section{Approximate segment integration, posterior prediction, and evaluation}
\label{sec:approx-predict}

For a nonconjugate segment model, \BayesBreak requires a numerical approximation to the segment marginal likelihood. Posterior prediction requires the posterior predictive distribution associated with the fitted observation model. This section states the error condition under which segmentwise approximation error can be propagated through the partition recursions and then defines the prediction and boundary-evaluation rules used in the empirical analysis.

\subsection{Propagation of segment marginal-likelihood error}
Let $\ell_{ij}=\log A^{(0)}_{ij}$ be a reference log segment marginal likelihood and let $\widehat\ell_{ij}$ be the value returned by a numerical method such as Laplace approximation, the Jaakkola--Jordan bound, expectation propagation, P\'olya--Gamma mean-field approximation, or adaptive quadrature.

\begin{assumption}[Uniform error bound over admissible segments]
\label{ass:block-error}
There exists $\varepsilon\ge0$ such that
\begin{equation}
 |\widehat\ell_{ij}-\ell_{ij}|\le\varepsilon
\label{eq:block-error}
\end{equation}
for every segment $(i,j]$ that can occur in a partition with $k\le k_{\max}$. The numerical analysis must identify the reference value and, when applicable, separate optimization residual, truncation or tail error, quadrature error, and discrepancy from the reference calculation.
\end{assumption}



\begin{theorem}[Stability of marginal likelihoods and posterior odds]
\label{thm:stability-paper}
Under Assumption~\ref{ass:block-error}, every $k$-segment forward or backward quantity satisfies
\begin{align}
e^{-k\varepsilon}L_{k,j}
&\le\widehat L_{k,j}\le e^{k\varepsilon}L_{k,j},\nonumber\\
e^{-k\varepsilon}R_{k,i}
&\le\widehat R_{k,i}\le e^{k\varepsilon}R_{k,i}.
\label{eq:message-sandwich}
\end{align}
Consequently, for $k,k'\le k_{\max}$,
\begin{equation}
\left|
\log\frac{\widehat p(k\mid y)}{\widehat p(k'\mid y)}
-\log\frac{p(k\mid y)}{p(k'\mid y)}
\right|\le(k+k')\varepsilon,
\label{eq:k-odds-stability}
\end{equation}
and, for two locations $h,h'$ of the same ordered boundary $t_p$ under fixed $k$,
\begin{equation}
\left|
\log\frac{\widehat p(t_p=h\mid y,k)}{\widehat p(t_p=h'\mid y,k)}
-\log\frac{p(t_p=h\mid y,k)}{p(t_p=h'\mid y,k)}
\right|\le2k\varepsilon.
\label{eq:boundary-odds-stability}
\end{equation}
\end{theorem}
\begin{proof}
\textbf{Step 1 (pathwise perturbation).}
For any admissible partition of $k$ blocks, the approximate-to-exact path-weight ratio is $\exp\{\sum_{q=1}^k\delta_{t_{q-1},t_q}\}$, with every summand bounded in absolute value by $\varepsilon$. Consequently that ratio lies between $e^{-k\varepsilon}$ and $e^{k\varepsilon}$.
\textbf{Step 2 (sum over positive paths).}
Multiply the exact weight of each path by the same lower and upper bounds and sum. Positivity preserves both inequalities. The prefix and suffix path sets contain exactly $k$ segments each, proving the two message sandwich bounds in Equation~\eqref{eq:message-sandwich}.
\textbf{Step 3 (count posterior odds).}
For distinct counts $k,k'$, the ratio $p(k\mid y)/p(k'\mid y)$ equals $[p(k)/C_k]L_{k,n}$ divided by $[p(k')/C_{k'}]L_{k',n}$. The count priors and partition normalizers are unchanged by local numerical approximation. The two summed-path errors therefore add in the log ratio, yielding the bound $(k+k')\varepsilon$.
\textbf{Step 4 (ordered boundary odds).}
The unnormalized weight for $t_p=h$ is $L_{p,h}R_{k-p,h}$; it is a sum of paths each with exactly $p+(k-p)=k$ blocks. Its log perturbation is bounded by $k\varepsilon$. Comparing this weight at $h$ and $h'$ doubles the maximum error, proving the $2k\varepsilon$ statement.
\textbf{Step 5 (scope).}
The argument compares only positive exact and approximate scores on the same admissible support. A bound on typical blocks, a per-method asymptotic formula, or a large observed maximum score discrepancy does not establish the uniform hypothesis.
\end{proof}

\begin{corollary}[Consequences for probabilities and modal locations]
\label{cor:stability-paper}
If finite unnormalized weights satisfy $|\log\widehat w_a-\log w_a|\le\eta$, then their normalized probabilities satisfy $e^{-2\eta}\le\widehat p_a/p_a\le e^{2\eta}$ and $\|\widehat p-p\|_{\mathrm{TV}}\le\min\{1,e^{2\eta}-1\}$. Moreover, a fixed-$k$ modal boundary separated from every alternative by exact log-odds margin $\Delta$ is unchanged whenever $\varepsilon<\Delta/(2k)$.
\end{corollary}
\begin{proof}
\textbf{Step 1 (bound the perturbed normalizer).}
Write $\widehat w_a=w_ae^{\delta_a}$ with $|\delta_a|\le\eta$. Let $p_a=w_a/\sum_bw_b$. Then $\sum_bp_be^{\delta_b}\in[e^{-\eta},e^\eta]$.
\textbf{Step 2 (normalized probability ratios and total variation).}
Since $\widehat p_a/p_a=e^{\delta_a}/\sum_bp_be^{\delta_b}$, this ratio lies in $[e^{-2\eta},e^{2\eta}]$. Therefore $|\widehat p_a-p_a|\le(e^{2\eta}-1)p_a$. Summing its positive part bounds total variation by $\min\{1,e^{2\eta}-1\}$.
\textbf{Step 3 (preserve the maximizing boundary).}
Suppose $h^\star$ exceeds every rival $h$ by exact log-odds margin at least $\Delta$. Theorem~\ref{thm:stability-paper} changes each pairwise log-odds difference by at most $2k\varepsilon$. If $2k\varepsilon<\Delta$, all perturbed differences remain positive, so $h^\star$ remains the unique modal boundary location.
\end{proof}

\begin{proofobligation}[Error control for individual numerical methods]
\label{obl:routine-certification-paper}
This paper does not establish a uniform convergence rate for Laplace approximation, the Jaakkola--Jordan bound, expectation propagation, P\'olya--Gamma mean-field approximation, or adaptive quadrature over all admissible segments. To apply Theorem~\ref{thm:stability-paper}, a nonconjugate analysis must establish Assumption~\ref{ass:block-error} under conditions appropriate to the observation family, parameter domain, segment support, initialization, convergence event, and tail behavior. Without that step, the dynamic program is exact for the supplied approximate score array, but the theorem does not compare the resulting posterior with the posterior under the intended segment model.
\end{proofobligation}

\subsection{Nonconjugate block calculations: what each approximation establishes}
\label{sec:nonconjugate-derivations}
For a canonical Binomial--logistic block with a Gaussian prior on its intercept, the exact local evidence is a one-dimensional integral. It provides a numerical reference against which deterministic approximations can be tested before the resulting segment scores enter the same partition recursion. A local approximation formula, its variational lower bound, and a rigorous numerical error bound are three different objects.

\begin{lemma}[Finite, unique logistic-block mode under a proper normal prior]
\label{lem:logistic-mode-paper}
Let $y_r\mid\theta\sim\operatorname{Binom}(m_r,\sigma(\theta))$ independently, with $m_r\ge0$, $\sigma(\theta)=(1+e^{-\theta})^{-1}$, and $\theta\sim\mathcal N(\nu,\rho^2)$ for finite $\rho^2>0$. Then the block log posterior is strictly concave and has a unique finite maximizer, including blocks containing only successes or only failures.
\end{lemma}
\begin{proof}
\textbf{Step 1 (posterior objective).} Apart from its Binomial coefficients, the log posterior is
\[
\Psi_B(\theta)=Y_B\theta-M_B\log(1+e^\theta)
-\frac{(\theta-\nu)^2}{2\rho^2}+\mathrm{constant},
\]
where $Y_B=\sum y_r$ and $M_B=\sum m_r$.
\textbf{Step 2 (strict curvature).} Direct differentiation gives
\[
\Psi_B''(\theta)=-M_B\sigma(\theta)\{1-\sigma(\theta)\}-\rho^{-2}<0
\]
for every finite $\theta$, even when $M_B=0$.
\textbf{Step 3 (existence, not just uniqueness).} The Binomial log-likelihood is bounded above by the sum of its log combinatorial coefficients. The negative quadratic prior term tends to $-\infty$ as $|\theta|\to\infty$. Hence $\Psi_B$ is coercive, attains a finite maximum by continuity, and strict concavity makes that maximizer unique. This explicit tail condition is why the conclusion does not follow from strict concavity alone under an arbitrary prior.
\end{proof}

\subsubsection{Laplace evaluation and its limitations}
Let $\widehat\theta_B$ be the mode and $H_B=-\Psi_B''(\widehat\theta_B)>0$. The first-order Laplace approximation is
\begin{equation}
\log\widehat A_B^{\rm Lap}=\Psi_B(\widehat\theta_B)
+\tfrac12\log(2\pi)-\tfrac12\log H_B.
\label{eq:laplace-paper}
\end{equation}
This is a local second-order approximation to the integral, \emph{not} a nonasymptotic upper bound on the absolute log-evidence error. A controlled asymptotic remainder would require derivatives beyond second order, concentration of the integral near the mode, and bounds on the tails outside that neighborhood. Numerical evaluation should use safeguarded Newton steps, check finite curvature and convergence, and fall back to adaptive one-dimensional integration on failure. The software's local curvature diagnostic does not by itself establish Assumption~\ref{ass:block-error}.

\subsubsection{Jaakkola--Jordan lower bound}
\begin{proposition}[Valid logistic tangent bound and block-evidence lower bound]
\label{prop:jj-paper}
For $\xi>0$ put $\lambda(\xi)=\tanh(\xi/2)/(4\xi)$ and $\sigma(x)=(1+e^{-x})^{-1}$. Then, for every real $x$,
\begin{equation}
\log\sigma(x)\ge\log\sigma(\xi)+\frac{x-\xi}{2}
-\lambda(\xi)(x^2-\xi^2).
\label{eq:jj-correct-paper}
\end{equation}
For a $\operatorname{Binom}(m_r,\sigma(\theta))$ observation, this quadratic lower bound has linear coefficient $y_r-m_r/2$ and quadratic coefficient $-m_r\lambda(\xi_r)$. Combining these terms with a Gaussian prior and integrating the resulting Gaussian kernel gives a deterministic \emph{lower} bound on $\log A^{(0)}_B$; this bound alone supplies no upper bound on the unknown error gap.
\end{proposition}
\begin{proof}
\textbf{Step 1 (rewrite the log sigmoid).} The identity $\log\sigma(x)=x/2-\log2-\log\cosh(x/2)$ separates a linear term from an even function of $x$.
\textbf{Step 2 (bound the even term).} Define $f(v)=\log\cosh(\sqrt v/2)$ for $v\ge0$. Its derivative for $v>0$ is $f'(v)=\tanh(\sqrt v/2)/(4\sqrt v)$, which decreases with $v$. Thus $f$ is concave and lies below its tangent at $v=\xi^2$: $f(x^2)\le f(\xi^2)+\lambda(\xi)(x^2-\xi^2)$.
\textbf{Step 3 (recover the inequality).} Substitute this tangent inequality into the expression for $\log\sigma(x)$, then use the identity at $x=\xi$; Equation~\eqref{eq:jj-correct-paper} follows, with equality at $x=\pm\xi$ in the even part.
\textbf{Step 4 (apply to Binomial data).} Write $\log p(y_r\mid\theta)=\log\binom{m_r}{y_r}+(y_r-m_r/2)\theta-m_r\log[2\cosh(\theta/2)]$. Apply the tangent bound to its even part and sum over observations. Exponentiating and integrating an everywhere-lower-bounded nonnegative integrand preserves the ordering. Thus $\log\widehat A_B^{\rm JJ}\le\log A_B$, but the value of $\log A_B-\log\widehat A_B^{\rm JJ}$ has not been upper-bounded by this argument.
\end{proof}

\subsubsection{P\'olya--Gamma augmentation and mean-field computation}
\begin{proposition}[Augmentation identity and Gaussian coordinate updates]
\label{prop:pg-paper}
For a Binomial likelihood without its $\theta$-independent coefficient, let $b_r=m_r$, $\kappa_r=y_r-m_r/2$ and denote the P\'olya--Gamma density by $p_{\rm PG}(\omega\mid b,0)$. Then
\begin{equation}
\frac{e^{y_r\theta}}{(1+e^\theta)^{m_r}}
=2^{-b_r}e^{\kappa_r\theta}
\int_0^\infty e^{-\omega\theta^2/2}p_{\rm PG}(\omega\mid b_r,0)\dd\omega.
\label{eq:pg-identity-paper}
\end{equation}
With $\theta\sim\mathcal N(\nu,\rho^2)$, the mean-field family $q(\theta)\prod_rq(\omega_r)$ admits coordinate updates
\begin{align}
q(\theta)&=\mathcal N(\mu_q,V_q),
& V_q^{-1}&=\rho^{-2}+\sum_r\mathbb E_q[\omega_r],\nonumber\\
\mu_q&=V_q\left(\nu/\rho^2+\sum_r\kappa_r\right),
& q(\omega_r)&={\rm PG}(b_r,c_r),
& c_r^2&=\mu_q^2+V_q,
\label{eq:pg-updates-paper}
\end{align}
where $\mathbb E_q[\omega_r]=b_r\tanh(c_r/2)/(2c_r)$ with its continuous limit at $c_r=0$. If an additional \emph{whole-likelihood} power $a_r$ is used, the identity requires $(b_r,\kappa_r)=(a_rm_r,a_r(y_r-m_r/2))$; multiplying an independent ${\rm PG}(m_r,0)$ variable by $a_r$ is not equivalent.
\end{proposition}
\begin{proof}
\textbf{Step 1 (factor logistic likelihood).} The identity $1+e^\theta=2e^{\theta/2}\cosh(\theta/2)$ converts the Binomial likelihood into $2^{-m_r}e^{\kappa_r\theta}\cosh(\theta/2)^{-m_r}$.
\textbf{Step 2 (introduce the PG mixture).} The P\'olya--Gamma Laplace-transform identity is $\cosh(\theta/2)^{-b}=\int e^{-\omega\theta^2/2}p_{\rm PG}(\omega\mid b,0)\dd\omega$ for $b>0$. Substitution proves Equation~\eqref{eq:pg-identity-paper} (with the $b=0$ contribution interpreted as one).
\textbf{Step 3 (update the Gaussian factor).} The augmented log joint is quadratic in $\theta$ conditional on the PG variables. Taking the expectation over $q(\omega)$ and adding the Gaussian prior gives precision $\rho^{-2}+\sum_r\mathbb E_q[\omega_r]$ and linear coefficient $\nu/\rho^2+\sum_r\kappa_r$, proving the Gaussian update.
\textbf{Step 4 (update the latent factors).} The coordinate log density for $q(\omega_r)$ is $\log p_{\rm PG}(\omega_r\mid b_r,0)-\omega_r\mathbb E_q[\theta^2]/2$ plus a constant. Exponential tilting gives ${\rm PG}(b_r,c_r)$, $c_r^2=\mathbb E_q[\theta^2]$, and the standard PG mean gives Equation~\eqref{eq:pg-updates-paper}. Each exact coordinate maximization cannot lower its \emph{properly evaluated} ELBO; it does not certify the gap to the exact evidence.
\textbf{Step 5 (external powers).} Raising the complete Binomial factor to $a_r$ changes both the hyperbolic-cosine exponent and linear coefficient, hence the PG shape and $\kappa_r$ as stated. A scaled PG random variable does not have this tilted law.
\end{proof}

For the intercept-only logistic--Gaussian block, the optimized Jaakkola--Jordan quadratic bound and the corresponding PG mean-field Gaussian updates can be algebraically equivalent. For the present numerical comparison, the \texttt{jj} and \texttt{pg\_vb} labels correspond to the same implementation and hence represent a single computational method. A distinct P\'olya--Gamma implementation would require separate numerical validation.

\subsubsection{Expectation propagation and numerical reference}
Expectation propagation approximates likelihood contributions with Gaussian sites and iterates cavity/tilted-distribution moment matching. A valid approximation to the \emph{segment marginal likelihood} must include the global Gaussian normalizer and the site/tilted normalizing corrections, not merely integrate a normalized approximate posterior (whose integral is one). The EP estimates are not used to establish integrated-likelihood accuracy in the absence of a separately validated normalizer, damping scheme, convergence criterion, and reference comparison. Adaptive one-dimensional quadrature, with a documented integration tolerance and tail or endpoint control, remains the numerical reference for low-dimensional blocks.


\subsection{Family-specific posterior prediction}
A fitted model records the observation family, the segment posterior state, and the coordinate range used for fitting. Prediction must use the posterior predictive distribution of that observation family. Coordinates outside the fitted range are rejected by default; an extrapolation rule is permitted only when it is stated as part of the predictive model before evaluation.



Consider a fitted segment $B$ with posterior conjugate hyperparameters $(\alpha_B,\beta_B)$. Let $(S_B^{\rm new},W_B^{\rm new},H_B^{\rm new})$ denote the exposure-aware sufficient summaries of new observations assigned to $B$.

\begin{proposition}[Posterior-predictive likelihood for a fixed segment]
\label{prop:predictive-paper}
Under the exponential-family formulation of Section~\ref{sec:model},
\begin{equation}
 p(y_B^{\rm new}\mid y_B)
 =e^{H_B^{\rm new}}
 \frac{Z(\alpha_B+S_B^{\rm new},\beta_B+W_B^{\rm new})}
      {Z(\alpha_B,\beta_B)}.
\label{eq:posterior-predictive}
\end{equation}
\end{proposition}
\begin{proof}
\textbf{Step 1 (begin with the fitted segment posterior).}
Conditional on the selected fitted segment, the training data update the conjugate hyperparameters from $(\alpha_0,\beta_0)$ to $(\alpha_B,\beta_B)$. This posterior, not the original segment prior, is the integrating measure for future observations.
\textbf{Step 2 (aggregate new sampling information).}
Under the same observation family and its declared descriptors, the new data likelihood has the kernel $\exp\{H_B^{\rm new}+\eta(\theta)^\top S_B^{\rm new}-W_B^{\rm new}A(\theta)\}$.
\textbf{Step 3 (integrate once).}
Multiplication by the training posterior produces the conjugate kernel with parameters $(\alpha_B+S_B^{\rm new},\beta_B+W_B^{\rm new})$. Integration yields its normalizer divided by $Z(\alpha_B,\beta_B)$, multiplied by $e^{H_B^{\rm new}}$, giving Equation~\eqref{eq:posterior-predictive}.
\textbf{Step 4 (conditioning).}
This density is posterior predictive conditional on the fitted segment and its assignment rule. It is not automatically averaged over alternative partitions, nor is it valid for a test family or coordinate support different from the declared predictive model.
\end{proof}

Prediction may condition on an exported joint-MAP partition or average over partitions and segment counts. The first calculation is conditional on a selected partition; the second integrates segmentation uncertainty. In either case, assignment of prediction coordinates to segments is part of the statistical model. Assigning every out-of-range coordinate to the final segment without an explicit extrapolation assumption is not a valid posterior-predictive calculation.

\subsection{Joint prediction for sequences, grouped units, and multivariate responses}
\label{sec:prediction-extensions}
The fixed-segment predictive formula in Proposition~\ref{prop:predictive-paper} supports several input structures, but \emph{joint} scoring must preserve dependence induced by a shared segment parameter. A product of separately integrated single-observation predictive densities is generally not the integrated density of their joint collection.

\begin{proposition}[Joint predictive integration within a fitted segment]
\label{prop:joint-predictive-paper}
Condition on a fitted partition $t$ and conjugate posterior $p(\theta_B\mid y_B)$ for each of its segments $B$. If new observations $D_B^{\rm new}$ assigned to $B$ are conditionally independent given $\theta_B$ and have sufficient summaries $(S_B^{\rm new},W_B^{\rm new},H_B^{\rm new})$, then
\begin{equation}
p(D_B^{\rm new}\mid y_B,t)
=e^{H_B^{\rm new}}\frac{Z(\alpha_B+S_B^{\rm new},\beta_B+W_B^{\rm new})}{Z(\alpha_B,\beta_B)}.
\label{eq:joint-predictive-block}
\end{equation}
The full fixed-partition joint predictive density is the product of these contributions over \emph{distinct segments}, not in general over individual observations within one segment.
\end{proposition}
\begin{proof}
\textbf{Step 1 (conditional likelihood).} Given $\theta_B$, independence gives $p(D_B^{\rm new}\mid\theta_B)=\prod_{r\in D_B^{\rm new}}p(y_r\mid\theta_B,d_r)$.
\textbf{Step 2 (aggregate once).} The sufficient-statistic representation collects all new observations into one exponential-family kernel with summaries $S_B^{\rm new},W_B^{\rm new},H_B^{\rm new}$.
\textbf{Step 3 (integrate once).} Multiply this kernel by the training posterior and integrate $\theta_B$, yielding Equation~\eqref{eq:joint-predictive-block}. It differs from integrating $\theta_B$ separately for each observation, which would spuriously treat new observations as having independent latent parameters.
\textbf{Step 4 (between-segment factorization).} Under the model assumption of conditionally independent segment parameters, the posterior factorizes across distinct segments given $t$. Therefore the joint predictive density factors over the distinct segments, completing the claim.
\end{proof}


\begin{corollary}[Closed-form predictive distributions in three conjugate families]
\label{cor:predictive-families-paper}
Conditional on a fitted segment and its conjugate posterior, the following are ordinary posterior-predictive sampling distributions:
\begin{enumerate}
\item For Gaussian $\mu\mid y_B\sim\mathcal N(\nu_B,V_B)$ and new observations with known variances $\sigma^2/w_r$, a vector of $m$ new observations in the same segment has joint law
\[
Y^{\rm new}\mid y_B,t\sim\mathcal N_m\!\left(\nu_B\mathbf1_m,\;V_B\mathbf1_m\mathbf1_m^\top+\operatorname{diag}(\sigma^2/w_1,\ldots,\sigma^2/w_m)\right).
\]
\item For Poisson $\lambda\mid y_B\sim\operatorname{Gamma}(a_B,b_B)$ and a new count at exposure $e>0$, the predictive law is negative Binomial with
\[
p(y^{\rm new}=z\mid y_B,e)=\frac{\Gamma(a_B+z)}{\Gamma(a_B)z!}\left(\frac{b_B}{b_B+e}\right)^{a_B}\left(\frac{e}{b_B+e}\right)^z.
\]
\item For Binomial $p\mid y_B\sim\operatorname{Beta}(a_B,b_B)$ and $m$ future trials, the predictive count law is Beta--Binomial with
\[
p(y^{\rm new}=z\mid y_B,m)=\binom mz\frac{B(a_B+z,b_B+m-z)}{B(a_B,b_B)}.
\]
\end{enumerate}
\end{corollary}
\begin{proof}
\textbf{Step 1 (Gaussian joint law).}
Write each new observation as $Y_r=\mu+\epsilon_r$ with independent $\epsilon_r\sim\mathcal N(0,\sigma^2/w_r)$ conditional on $\mu$, and with $\mu\sim\mathcal N(\nu_B,V_B)$ from the training posterior. Gaussian addition gives mean $\nu_B\mathbf1_m$ and covariance $V_B\mathbf1_m\mathbf1_m^\top+\operatorname{diag}(\sigma^2/w_r)$. The rank-one term shows why multiplying separate Gaussian marginals would be incorrect for joint scoring.
\textbf{Step 2 (Poisson--Gamma mixture).}
Integrate $e^{-e\lambda}(e\lambda)^z/z!$ against the Gamma posterior density. The Gamma integral has shape $a_B+z$ and rate $b_B+e$, which simplifies to the displayed negative-Binomial probability mass function.
\textbf{Step 3 (Beta--Binomial mixture).}
Integrate $\binom mzp^z(1-p)^{m-z}$ against $\operatorname{Beta}(a_B,b_B)$. Combining exponents produces $B(a_B+z,b_B+m-z)/B(a_B,b_B)$, establishing the third expression.
\end{proof}

\paragraph{Pointwise sequences (Case A).}
New pointwise observations carry coordinates, outcomes, and the family-specific descriptors. Under an exported partition, coordinates select segments and the data are aggregated \emph{within} those segments before applying Equation~\eqref{eq:joint-predictive-block}. An out-of-range coordinate requires a declared extrapolation model; a coordinate gap is not an exposure by default.

\paragraph{Set-valued units (Case B).}
A unit may contain multiple pairs $(x_{ur},y_{ur})$ on a physical support interval. Its segment-specific internal observations can be scored jointly using the same sufficient-statistic aggregation. If several units in a single dataset share segment parameters, the dataset score aggregates \emph{across all such units before integrating}. Multiplying separately marginalized unit likelihoods is justified only under an additional model in which the units have independent latent segment parameters; it is otherwise a composite score.

\paragraph{Vector-valued responses (Case C).}
When response coordinates are conditionally independent given distinct component-specific segment parameters, the joint segment evidence and predictive density factor across response dimensions. For a genuinely multivariate exponential family with shared or correlated parameters, the joint sufficient statistics and conjugate normalizer must be used instead. The global ordered-partition recursion depends only on the resulting joint segment score.

\paragraph{Group membership and segmentation averaging.}
If $M_g$ defines a \emph{normalized} predictive data density and $\pi_g$ is a prior group probability, then $p(g\mid D^{\rm new})\propto\pi_gp(D^{\rm new}\mid M_g)$ is a Bayesian group posterior. Scores from the finite latent-group MM criterion in Equation~\eqref{eq:latent-group-objective}, which are not asserted to be normalized data densities, do not automatically admit this interpretation. Likewise, true segmentation-averaged prediction sums joint predictive contributions over $p(k,t\mid y)$; prediction conditional on a chosen joint-MAP partition is a different object.

\paragraph{Calibration.}
For discrete outcomes, randomize the probability integral transform within the probability mass at the observed outcome. For continuous outcomes, PIT values are marginally uniform under a calibrated predictive law, but they are generally dependent within a sequence. Report sequence-level calibration or simulation-based reference distributions rather than applying an independent-uniform goodness-of-fit reference blindly. Held-out predictive scores need not rank candidate $k$ in the same order as training posterior count probabilities.


\subsection{Boundary metrics}
\begin{definition}[One-to-one tolerance matching]
\label{def:boundary-matching-paper}
For estimated boundaries $\widehat{\mathcal B}$, reference boundaries $\mathcal B^\star$, and tolerance $\tau$, form the bipartite graph containing an edge $(\widehat b,b^\star)$ whenever $|\widehat b-b^\star|\le\tau$. Choose a matching of maximum cardinality; among maximum-cardinality matchings, choose one of minimum total absolute distance under a deterministic tie rule. Precision, recall, and $F_1$ use the resulting number of matches. No boundary is used more than once.
\end{definition}

\begin{definition}[Matched location error]
\label{def:matched-mae-paper}
For the matching $\mathcal M_\tau$,
\[
\operatorname{MAE}_\tau
=|\mathcal M_\tau|^{-1}
 \sum_{(\widehat b,b^\star)\in\mathcal M_\tau}|\widehat b-b^\star|
\]
when $|\mathcal M_\tau|>0$. When no pair is matched, the value is reported as \NA, not as zero.
\end{definition}

\begin{definition}[Admissible held-out predictive score]
\label{def:heldout-score-paper}
A held-out log score is used as predictive evidence only when the held-out observations were excluded from fitting, the posterior predictive distribution of the fitted observation family was used, and every prediction coordinate lay in the fitted support or followed a pre-specified extrapolation model. A violation of any condition invalidates the number for that predictive interpretation without changing the fact that the computation was executed.
\end{definition}

\paragraph{Predictive model checking.}
A posterior-predictive score is meaningful only for the observation law that
was fitted and for new coordinates within its stated support. An earlier
methylation calculation used a Gaussian predictive density for a fitted
Beta-response model and extrapolated outside the fitted range. That value is
not included in the predictive results. The admissible Beta-response
held-out analysis is reported in Section~\ref{sec:results}.

\paragraph{Interpretation of approximate local scores.}
The forward--backward and max--sum recursions are exact for the block-score array supplied to them. When the array is obtained by numerical approximation, this algorithmic exactness does not establish exactness with respect to the underlying sampling model. A Jaakkola--Jordan tangent gap averaged under a fitted Gaussian variational distribution is not an upper bound on the log-evidence error: the difference in log normalizing constants also depends on the changed posterior measure. Likewise, a truncated Laplace expansion is not a finite-sample error certificate unless its omitted remainder and tails have been bounded. Comparisons with a converged numerical-integration reference are necessary before interpreting either diagnostic as evidence of approximation accuracy.
\paragraph{Interpretation of nonconjugate method labels.}
In the implementation studied here, the options \texttt{jj} and \texttt{pg\_vb} both invoke the same Jaakkola--Jordan numerical routine. These labels therefore do not constitute two independent approximation methods. A distinct P\'olya--Gamma variational implementation, if included later, requires a separately specified augmented model, objective, update equations, and numerical verification.
